\section{Key-Recovery Security of Classical Schnorr} \label{sec-lf-kr}

Our first result concerns key recovery rather than unforgeability. It shows that
an adaptive signing oracle gives no additional help in recovering the Schnorr
secret key for a particular standard-model hash family. The proof uses a lossy
mode with an efficiently enumerable image to simulate signatures without the
secret key. The real scheme uses the injective mode, which is indistinguishable
from the lossy mode. We later explain why this technique stops at key recovery:
if the small image used by the simulator were publicly enumerable in the real
scheme, the same simulation would be a forgery attack.

\subsection{Key Recovery under Chosen-Message Attack}

For a signature scheme $\SigScheme=(\SigKeys,\SigSign,\SigVf)$, the key-recovery
game gives the adversary the verification key and an adaptive signing oracle,
and the adversary wins by returning the signing key:
\begin{center}
\fbox{
\begin{minipage}{2.8in}
\ExperimentHeader{Game $\gameBKCMA{\SigScheme}$}

\begin{oracle}{$\Initialize$}
\item $(\vkey,\skey) \getsr \SigKeys$
\item Return $\vkey$
\end{oracle}
\ExptSepSpace
\begin{oracleC}{$\SignOracle(m)$}{$m \in \MsgSp$}
\item $\sigma\getsr\SigSign(\skey,m)$
\item Return $\sigma$
\end{oracleC}
\ExptSepSpace
\begin{oracle}{$\Finalize(\skey')$}
\item Return $(\skey'=\skey)$
\end{oracle}
\end{minipage}
}
\end{center}
The BK-CMA advantage of $\advA$ is
$$
\AdvBKCMA{\SigScheme}{\advA}
:=
\Pr\big[\gameBKCMA{\SigScheme}(\advA)\Rightarrow1\big].
$$
For Schnorr signatures the signing key is $(x,\hk)$, and $\hk$ is included in
the public verification key. Hence an adversary that returns $x$ can return
$(x,\hk)$. In the other direction, when an adversary returns $(x',\hk')$, output
$x'$ if $\hk'=\hk$ and output $\bot$ otherwise. This transformed adversary
recovers $x$ exactly when the original adversary returns the full signing key.
Thus the two formulations have the same advantage. We use the
exponent-recovery formulation below.

\subsection{The Hash Family} \label{sec-lf-constr}

Let $\GG$ be a cyclic group of prime order $p$ with generator $g$. Let $\TLF$ be
a tunable lossy function with algorithms $\TLFKg,\TLFEv,\TLFEnum$. Its domain is
$\bits^{\ellin}$ and its range is $\Ygrave$. Recall from Section~\ref{sec-lf}
that $M=2^{\ellin}$, and fix an injective encoding
$$
\enc{\cdot}\colon\GG\times\bits^{\msglen}\longrightarrow\bits^{\ellin}.
$$

Schnorr requires its challenge to lie in $\Zp$, whereas the range $\Ygrave$ of
the injective mode must be large enough to contain an injective image of all
pairs $(R,m)$. We therefore use an efficient keyed output-map family
$$
\mathsf{Map}=(\mathsf{MapKg},\mathsf{MapEv}).
$$
Algorithm $\mathsf{MapKg}$ on input $\usecp$ outputs a key $k$, and the
deterministic algorithm $\mathsf{MapEv}$ on inputs $k$ and $y\in\Ygrave$ outputs
an element of $\Zp$. No security property of $\mathsf{Map}$ is required for the
key-recovery result.

\begin{construction} \label{constr-lf}
Fix an image parameter $r\in[M]$. The keyed hash family
$\HF^{\mathsf{tlf}}_r=(\HFKg,\HFEv)$ is defined as follows. Algorithm $\HFKg$ on
input $X\in\GG$ computes
$$
(K,\mathsl{td})\getsr\TLFKg(\usecp,r,\mathsf{inj})
\qquad\text{and}\qquad
k\getsr\mathsf{MapKg}(\usecp),
$$
and returns $\hk=(K,k)$. Algorithm $\HFEv$ on inputs $\hk=(K,k)$ and
$(R,m)\in\GG\times\bits^{\msglen}$ returns
$$
\mathsf{MapEv}\bigl(k,\TLFEv(K,\enc{R,m})\bigr).
$$
\end{construction}

Key generation ignores $X$, and lossy keys are used only in the proof. The
output map has a syntactic role in the construction: it maps the possibly
structured range $\Ygrave$ to the scalar challenge required by Schnorr. This
compression cannot in general be absorbed into the TLF while retaining the
injectivity required by Definition~\ref{def-lf}. Indeed, the encoding domain
contains $p2^{\msglen}$ pairs, while $\Zp$ contains only $p$ elements. Thus it is
the inner TLF generated in injective mode that is injective, except with probability
$\nu_{\mathsf{inj}}(\secp,r)$. The final composite hash is not claimed to be
injective.

\subsection{Exact Rejection Sampling} \label{sec-bk-sampling}

The proof rests on the following elementary fact. It is independent of lossy
functions and applies to any fixed Schnorr challenge function whose values, for
the message being signed, lie in a known finite set.

\begin{lemma}[Exact Schnorr rejection sampling] \label{lem-sim}
Fix $x\in\Zp$ and $X=g^x$. Let
$H\colon\GG\times\bits^{\msglen}\to\Zp$ be a function, let
$m\in\bits^{\msglen}$, and let $C\subseteq\Zp$ be a nonempty finite set such
that
$$
\{H(R,m):R\in\GG\}\subseteq C.
$$
Consider one trial that samples $c\getsr C$ and $s\getsr\Zp$, sets
$R\gets g^sX^{-c}$, and accepts exactly when $H(R,m)=c$. The trial accepts with
probability $1/|C|$. Conditioned on acceptance, $(R,s)$ has exactly the
distribution of an honest Schnorr signature on $m$. Consequently, for every
integer $N\ge1$, $N$ independent trials all fail with probability at most
$e^{-N/|C|}$; in particular, $\secp|C|$ trials all fail with probability at
most $e^{-\secp}$.
\end{lemma}

\begin{proof}
Call $(c,s)\in C\times\Zp$ good if
$$
H\bigl(g^sX^{-c},m\bigr)=c.
$$
For each $R\in\GG$, let $\rho_R\in\Zp$ be the unique scalar such that
$R=g^{\rho_R}$,
and define
$$
c_R:=H(R,m)
\qquad\text{and}\qquad
s_R:=\rho_R+c_Rx\pmod p.
$$
The covering condition gives $c_R\in C$, and
$g^{s_R}X^{-c_R}=R$. Thus $(c_R,s_R)$ is good. Conversely, if $(c,s)$ is good
and $R=g^sX^{-c}$, then $c=c_R$ and the same equation gives $s=s_R$. Hence the
good pairs are exactly
$$
\{(c_R,s_R):R\in\GG\},
$$
and the map $(c,s)\mapsto g^sX^{-c}$ is a bijection from the good pairs to
$\GG$.

There are therefore exactly $p$ good pairs among the $|C|p$ equally likely
pairs in $C\times\Zp$, so one trial accepts with probability $1/|C|$.
Conditioned on acceptance, the good pair is uniform, and the bijection makes
$R$ uniform in $\GG$. For this $R$, the response is
$s=\rho_R+H(R,m)x\bmod p$, exactly as in honest Schnorr signing. Finally,
$(1-1/|C|)^N\le e^{-N/|C|}$.
\end{proof}

\subsection{Main Result} \label{sec-lf-thm}

\paragraph*{Proof idea.}
Fix the image parameter $r$. We first switch the inner function from its
injective mode to its indistinguishable lossy mode. The reduction enumerates a
set covering the lossy image and applies the output map to obtain a set
$C\subseteq\Zp$ covering all possible Schnorr challenges. It then answers every
adaptive signing query using Lemma~\ref{lem-sim}. The answers have exactly the
honest distribution, apart from the stated enumeration and truncation failure
probabilities, and no signing query uses $x$. A key-recovery adversary can
therefore be run on a DL challenge.

\begin{theorem} \label{thm-bk}
Let $\advA$ be a PPT adversary against BK-CMA of
$\SchScheme[\GG,\HF^{\mathsf{tlf}}_r]$ that runs in time $t_{\advA}$ and makes
at most $q_s$ signing queries. For every $r\in[M]$, there are a distinguisher
$\advD$ and a DL algorithm $\advB$ such that
$$
\AdvBKCMA{\SchScheme[\GG,\HF^{\mathsf{tlf}}_r]}{\advA}
\le
\AdvTLF{\TLF,r}{\advD}
+\AdvDL{\GG,g}{\advB}
+\nu_{\mathsf{enum}}(\secp,r)
+q_s e^{-\secp}.
$$
Adversary $\advD$ runs in time $t_{\advA}+q_s\poly(\secp)$, and $\advB$ runs
in time
$$
t_{\advA}+\tau(\secp,r)+(q_s+1)\secp r\poly(\secp).
$$
\end{theorem}

The quantitative reduction uses mode indistinguishability and enumeration of
the lossy image, together with DL in $\GG$. It does not use injectivity of the
injective mode or any security property of the output map. The construction
nevertheless generates its TLF key in injective mode; injectivity itself is
not used in the reduction. The theorem holds for every efficient output map.

The games are in Figure~\ref{fig-bk-games}. In the figure, write
$$
H(R,m)=\mathsf{MapEv}\bigl(k,\TLFEv(K,\enc{R,m})\bigr).
$$
Game $\Gm_0$ is the BK-CMA game. Game $\Gm_1$ changes the tunable lossy function
to lossy mode at the same value $r$. Game $\Gm_2$ uses the enumeration
information retained by the reduction to enumerate the lossy image once, and
replaces honest signing by rejection sampling. If $C$ is empty, which can occur
only within the enumeration-failure event, the oracle returns $\bot$.

\FloatBarrier
\begin{figure}[t]
\twoCols{0.42}{0.50}{
\ExperimentHeader{Game $\mathsf{G}[\mathsf{mode},b]$}

\begin{oracle}{$\Initialize$}
\item $x\getsr\Zp\;;\;X\gets g^x$
\item $k\getsr\mathsf{MapKg}(\usecp)$
\item $(K,\mathsl{td})\getsr\TLFKg(\usecp,r,\mathsf{mode})$
\item If $b=1$:
\item \quad $S\getsr\TLFEnum(\mathsl{td})$
\item \quad $C\gets\{\mathsf{MapEv}(k,y):y\in S\}$
\item $\hk\gets(K,k)$; Return $(X,\hk)$
\end{oracle}
\ExptSepSpace
\begin{oracle}{$\Finalize(x')$}
\item Return $(x'=x)$
\end{oracle}
}
{
\begin{oracle}{$\SignOracle(m)$}
\item If $b=0$:
\item \quad $\rho\getsr\Zp\;;\;R\gets g^{\rho}$
\item \quad $c\gets H(R,m)\;;\;s\gets(\rho+cx)\bmod p$
\item \quad Return $(R,s)$
\item If $C=\emptyset$: Return $\bot$
\item Repeat $\secp|C|$ times:
\item \quad $s\getsr\Zp\;;\;c\getsr C$
\item \quad $R\gets g^sX^{-c}$
\item \quad If $H(R,m)=c$: Return $(R,s)$
\item Return $\bot$
\end{oracle}
}
\vspace{-12pt}
\caption{Games for the proof of Theorem~\ref{thm-bk}. We set
$\Gm_0=\mathsf{G}[\mathsf{inj},0]$,
$\Gm_1=\mathsf{G}[\mathsf{lossy},0]$, and
$\Gm_2=\mathsf{G}[\mathsf{lossy},1]$.}\label{fig-bk-games}
\vspace{-3pt}\hrulefill
\vspace{-1ex}
\end{figure}
\FloatBarrier

\begin{proof}[of Theorem~\ref{thm-bk}]
\heading{From $\Gm_0$ to $\Gm_1$.}
The distinguisher $\advD$, on input $K$, samples $x\getsr\Zp$ and
$k\getsr\mathsf{MapKg}(\usecp)$. It sets $\hk=(K,k)$, runs $\advA$ on $(g^x,\hk)$, answers
signing queries honestly, and returns 1 exactly when $\advA$ outputs $x$. The
two input distributions for $\advD$ give games $\Gm_0$ and $\Gm_1$. The
enumeration information generated together with $K$ is hidden and unused in
both games, so $\advD$ needs only the public key supplied by the TLF game. Hence
$$
|\Pr[\Gm_0\Rightarrow1]-\Pr[\Gm_1\Rightarrow1]|
\le\AdvTLF{\TLF,r}{\advD}.
$$

\heading{From $\Gm_1$ to $\Gm_2$.}
Next condition on successful enumeration in $\Gm_2$. The simulator computes
$S\getsr\TLFEnum(\mathsl{td})$ and then
$C=\{\mathsf{MapEv}(k,y):y\in S\}$. For every message $m$,
$$
\{H(R,m):R\in\GG\}\subseteq C,
$$
because the enumerated set covers the image of $\TLFEv(K,\cdot)$ on its entire
domain. Lemma~\ref{lem-sim} therefore says that each signing query has exactly
the distribution in $\Gm_1$ unless its rejection loop fails. Hence
$$
|\Pr[\Gm_1\Rightarrow1]-\Pr[\Gm_2\Rightarrow1]|
\le\nu_{\mathsf{enum}}(\secp,r)+q_se^{-\secp}.
$$

\heading{From $\Gm_2$ to DL.}
Finally, algorithm $\advB$, on input a DL challenge $h$, sets $X\gets h$ and samples
$$
(K,\mathsl{td})\getsr\TLFKg(\usecp,r,\mathsf{lossy})
\qquad\text{and}\qquad
k\getsr\mathsf{MapKg}(\usecp),
$$
sets $\hk=(K,k)$, and runs $\advA$ on $(X,\hk)$. Since $\advB$ generated the lossy key, it retains
$\mathsl{td}$ and can compute the set $C$ used by $\Gm_2$. It answers signing
queries as in $\Gm_2$ and returns the value output by $\advA$. This perfectly
simulates $\Gm_2$, so
$$
\Pr[\Gm_2\Rightarrow1]=\AdvDL{\GG,g}{\advB}.
$$
The lossy image is enumerated once. Constructing $C$ takes
$\tau(\secp,r)+r\poly(\secp)$ time, and each signing query uses at most
$\secp r$ trials. This gives the claimed running times and completes the proof.
\end{proof}

\subsection{Why Lossiness Stops at Key Recovery}

The rejection-sampling loop uses only the verification key once a covering set
for the relevant challenge values is known. It therefore becomes a forgery
attack whenever such a small set is publicly computable. The following
proposition makes this precise and does not depend on how the hash function is
built.

\begin{proposition} \label{prop-lossy-forge}
Let $\HF=(\HFKg,\HFEv)$ be a keyed hash family for $\GG$. Suppose there is an
algorithm $\mathsf{Cover}$ that, on input a public verification key
$(X,\hk)$ in the support of Schnorr key generation, runs in time $\tau$ and
returns a message $m^*\in\bits^{\msglen}$ and a nonempty set
$C\subseteq\Zp$ such that $|C|\le r$ and
$$
C\supseteq\{\HFEv(\hk,(R,m^*)):R\in\GG\}.
$$
Then there is an adversary that runs in time
$\tau+\secp r\poly(\secp)$, makes no signing queries, and wins the EUF-CMA game
against $\SchScheme[\GG,\HF]$ with probability at least $1-e^{-\secp}$.
\end{proposition}

\begin{proof}
On input $(X,\hk)$, the adversary runs
$(m^*,C)\gets\mathsf{Cover}(X,\hk)$. It then repeats the rejection-sampling loop
of Figure~\ref{fig-bk-games} on $m^*$, using the public function
$H(R,m)=\HFEv(\hk,(R,m))$ and the set $C$. Lemma~\ref{lem-sim} says that each
trial succeeds with probability $1/|C|$. Thus all $\secp|C|$ trials fail with
probability at most $e^{-\secp}$. Any output is a valid signature on $m^*$, and
the adversary made no signing query.
\end{proof}

The qualification ``publicly computable'' is essential: the proposition needs
an efficient way to obtain a covering set, not merely an information-theoretic
bound on the image. The abstract TLF definition allows the reduction to retain
auxiliary enumeration information generated together with the lossy key. The
DDH- and LWE-based constructions of Propositions~\ref{prop-tlf-ddh}
and~\ref{prop-tlf-lwe} both use such information. The deployed key is instead
generated in injective mode, and the corresponding small cover is unavailable
to the public-cover attack above.

\paragraph*{Why not deploy the lossy mode?}
One could do so. More generally, the proof works directly for any fixed
challenge function for which a polynomial-size covering set can be efficiently
computed: the mode switch is then unnecessary, and the rejection sampler
reduces key recovery directly to DL. A constant challenge function is the
extreme example. It can still be BK-CMA secure, because seeing valid signatures
need not reveal $x$, but it is trivially forgeable: if its value is $c$, choose
$s\getsr\Zp$ and return $(g^sX^{-c},s)$ on any message.

This observation is why a deliberately small-image real hash says little about
the behavior of Schnorr with a practical challenge hash. Such a hash is intended
to spread its outputs throughout the challenge space, not to have a publicly
computable polynomial- or quasi-polynomial-size cover.

The point of Construction~\ref{constr-lf} is narrower. Its deployed inner
function is generated in the injective mode, while only the indistinguishable
lossy mode used in the proof has the small enumerable image. The composite
still has the usual Schnorr domain and range: it takes $(R,m)$ and returns a
scalar in $\Zp$. The
theorem permits any efficient output map, including degenerate choices, and
therefore proves key recovery only. It does not prove that SHA-3 has a lossy
mode, or that plain SHA-3 satisfies the theorem. It shows instead that the
small-image simulation need not be a visible property of the deployed inner
function. Lossiness alone cannot establish unforgeability, which is why the
results that follow use different tools.

\subsection{Assumption Tradeoffs}

Theorem~\ref{thm-bk} gives the tradeoff directly in terms of the TLF. For any
chosen $r=r(\secp)$, the mode distinguisher runs in time
$t_{\advA}+q_s\poly(\secp)$ and never enumerates the lossy image. The DL
algorithm instead runs in time
$$
t_{\advA}+\tau(\secp,r)+(q_s+1)\secp r\poly(\secp).
$$
Thus the TLF switch need only be secure against polynomial-time distinguishers,
even when the enumeration and the resulting DL reduction are
superpolynomial.

\begin{corollary}[Fixed-image tradeoff] \label{cor-bk-fixed}
Fix any efficiently represented image bound $r=r(\secp)$. Suppose every PPT
distinguisher has negligible TLF advantage at $r$, the enumeration failure
$\nu_{\mathsf{enum}}(\secp,r)$ is negligible, and, for every polynomial $P$,
every algorithm running for at most
$$
\tau(\secp,r)+rP(\secp)
$$
time has negligible DL advantage in $\GG$. Then
$\SchScheme[\GG,\HF^{\mathsf{tlf}}_r]$ is BK-CMA secure for every efficient
output-map family $\mathsf{Map}$.
\end{corollary}

\begin{proof}
Apply Theorem~\ref{thm-bk}. Its TLF distinguisher is PPT, while its DL
algorithm falls within the assumed running time after absorbing the
adversary's polynomial running time. The TLF advantage, DL advantage, and
enumeration-failure probability are negligible, as is $q_se^{-\secp}$ for
polynomial $q_s$.
\end{proof}

In particular, if $\tau(\secp,r)\le r\poly(\secp)$ and
$r=2^{\secp^\theta}$ for a constant $0<\theta<1$, the DL side requires hardness
against $2^{O(\secp^\theta)}$-time algorithms, while the TLF distinguisher
remains polynomial time. A quasi-polynomial $r$ gives the analogous
quasi-polynomial DL requirement. Propositions~\ref{prop-tlf-ddh}
and~\ref{prop-tlf-lwe} give DDH- and LWE-based ways to realize the TLF profile
at the chosen value of $r$. Under the formal definition of subexponential
security in Section~\ref{sec-prelims}, DL security with any exponent
$\delta_{\mathsf{DL}}>\theta$ absorbs the constant hidden in the displayed
$O(\secp^\theta)$ running time.

There is a stronger endpoint in which the DL reduction remains polynomial.
For a TLF satisfying Definition~\ref{def-tlf-strong}, the injective-key
distribution does not depend on the image parameter. Since the output-map key
is also sampled independently of $r$, the distribution of
$\HF^{\mathsf{tlf}}_r$ is the same for every $r$; we write
$\HF^{\mathsf{tlf}}$ for this family.

\begin{corollary} \label{cor-bk}
Suppose $\TLF$ is strongly efficiently enumerable and DL is hard against
PPT algorithms in $\GG$. Then
$\SchScheme[\GG,\HF^{\mathsf{tlf}}]$ is BK-CMA secure for every efficient
output-map family $\mathsf{Map}$.
\end{corollary}

\begin{proof}
Suppose a PPT adversary $\advA$ has nonnegligible BK-CMA advantage. Fix an
inverse polynomial $\epsilon$ such that this advantage is at least $4\epsilon$
for infinitely many security parameters. The distinguisher $\advD$ of
Theorem~\ref{thm-bk} is PPT, and its code is independent of $r$: it receives
$K$ and never uses the image bound. Let $q$ be the polynomial given by
Definition~\ref{def-tlf-strong} for a polynomial time bound on $\advD$ and for
$\epsilon$, and set
$r(\secp)=\lceil q(\secp)\rceil$. Since $q$ is polynomial and
$M\ge2^\secp$, we have $r(\secp)\le M$ for all sufficiently large $\secp$.
Then
$\AdvTLF{\TLF,r}{\advD}\le\epsilon$. The injective-mode distribution does not
depend on this choice of $r$, so the left side of the theorem is still the
advantage against the honest scheme.

For this fixed polynomial $r$, algorithm $\advB$ is PPT. Its DL advantage is
negligible. The enumeration-failure term is also negligible.
The remaining term $q_se^{-\secp}$ is negligible as well. For all sufficiently
large parameters the right side is less than $4\epsilon$, a contradiction.
\end{proof}

\paragraph*{Instantiation and comparison.}
By Proposition~\ref{prop-tlf-ddh}, Corollary~\ref{cor-bk} assumes ordinary DL
in the Schnorr group and exponential DDH, or constant-$k$ exponential
$k$-linear, in auxiliary groups. The cascade is what permits the proof to
choose a polynomial image bound after fixing the attack; the fixed-image
Corollary~\ref{cor-bk-fixed} does not require that stronger quantifier order.

Paillier and Vergnaud~\cite{AC:PaiVer05} prove BK-CMA security for Schnorr with
any hash function under OMDL. Their result covers any hash function, while ours
covers a particular keyed family. Our group assumption is ordinary DL, and the
DDH-cascade instantiation uses exponential DDH for the TLF. OMDL is
interactive; the assumptions used by our BK-CMA reduction are non-interactive.
